\section{Method}
\subsection{Problem formulation and cost}
Let source image $I_i$ belong to crop domain $d_i$ and occupy valid pixel domain $\Omega_i$. A tiling operator $\mathcal{T}_{s,\alpha}$ produces candidate tiles $\mathcal{T}_i$ with clipped source-coordinate footprints $\Omega_t\subseteq\Omega_i$ and centers $x_t$. Source images are partitioned before tiling. The detector observes every candidate tile; at most $K$ tiles receive VLM inference. The system predicts one record for the complete source image and separately records the union of VLM-inspected footprints. An uninspected region is therefore not treated as an observed negative.

Table~\ref{tab:schema} fixes the roles of all output fields. The four evaluated task fields are annotation-derived inventory surrogates, not independent agronomic labels. Free text, self-reported confidence, and the accept/review decision are audit or system fields and are never scored as biological ground truth.

\begin{table*}[t]
\centering
\wacvtablefont
\caption{Frozen source-record specification. Thresholds and class support are fixed from training data before test evaluation. Annotated coverage is box-derived and must not be described as canopy segmentation.}
\label{tab:schema}
\setlength{\tabcolsep}{4.5pt}
\begin{tabular}{lllll}
\toprule
Field & Allowed value or type & Reference source & Task metric & Gate input \\
\midrule
Plant presence & yes/no & any verified source box & macro-F1 & yes \\
Count & nonnegative estimate and frozen interval & unique source-coordinate boxes & MAE; interval F1 & yes \\
Density class & empty/sparse/moderate/dense & frozen count/valid-area rule & macro-F1 & yes \\
Annotated coverage & none/low/medium/high & clipped box union/valid area & macro-F1; error & yes \\
Visible evidence & text with tile provenance & no independent reference & audit only & no \\
Model confidence & $[0,1]$, self-reported & no independent reference & calibration diagnostic & candidate \\
Decision & accept/review & gate output & risk--coverage & output \\
\bottomrule
\end{tabular}
\end{table*}

For candidate-tile detector cost $C_D$, selected-tile VLM cost $C_V$, and source reduction cost $C_R$, the measured source cost is
\begin{equation}
C(i)=|\mathcal{T}_i|C_D+|\mathcal{S}_i|C_V+C_R,
\qquad |\mathcal{S}_i|\leq K.
\label{eq:cost}
\end{equation}
Selection controls the expensive second term. Rejection occurs after generation and is evaluated as risk control rather than compute saving.

\subsection{Detector evidence and controlled baselines}
A detector trained only on training sources returns source-coordinate detections $\mathcal{D}_t$, their count $n_t=|\mathcal{D}_t|$, and mean retained confidence $c_t$. Confidence thresholding, non-maximum suppression, valid-footprint clipping, and padding removal precede scoring; $c_t=0$ when no detection remains. After this processing, the ranking score is
\begin{equation}
q_t=c_t+\frac{\min(n_t,n_{\max})}{n_{\max}}.
\label{eq:evidence}
\end{equation}
The clipping value $n_{\max}$ and all detector post-processing parameters are chosen on validation sources; crop-stratified saturation rates are reported. Ties are resolved by source-coordinate row, column, then stable tile identifier. Confidence-only, count-only, normalized combined, and spatially deduplicated scores isolate the contribution of each term. The combined score is treated as a heuristic unless a validation comparison supports a stronger interpretation. A detector-only baseline maps fused source-coordinate boxes directly to all derivable output fields and never invokes a VLM.

Training-only hard-negative mining adds reviewed unmatched detector responses from empty training regions. Naturally empty candidates, boundary-aware empty candidates, mined false-positive crops, and verified negative evaluation units are tracked as distinct sets. No validation or test annotation is used to create training negatives. The mining detector, review procedure, and matched-exposure control are reported separately so that any gain is not attributed to evidence allocation.

\subsection{Fixed-budget evidence allocation}
For source $i$, let $K_i=\min(K,|\mathcal{T}_i|)$. The exploit--explore policy assigns $K_e=\min(|\mathcal{T}_i|,\operatorname{round}(\rho K_i))$ slots to exploration and $K_x=K_i-K_e$ to exploitation:
\begin{align}
\mathcal{X}_i &= \operatorname{TopK}_{K_x}(\mathcal{T}_i;q), \\
\mathcal{E}_i &\sim \operatorname{Sample}_{K_e}(\mathcal{T}_i\setminus\mathcal{X}_i), \\
\mathcal{S}_i &= \mathcal{X}_i\cup\mathcal{E}_i.
\label{eq:selector}
\end{align}
If the remaining exploration pool is smaller than $K_e$, every remaining tile is used and the unused slots are returned to exploitation. Sampling uses a seed derived from the study seed, source identifier, policy, budget, and repetition, making the result invariant to manifest order. Stochastic policies are evaluated over repeated seeds; reproducibility from one seed is not treated as robustness. Candidate manifests are shared by all policies, and whether budget sets are nested is declared before evaluation.

Equal-budget controls include random sampling, confidence top-$K$, count top-$K$, combined-score top-$K$, and two geometry-aware policies. Spatial selection applies farthest-point sampling to normalized tile centers $x_t$, while diversity-aware exploration selects remaining tiles using distance from the exploited set and detector-score uncertainty. Ground-truth routing and shuffled detector scores are diagnostic controls only. The method does not assume that exploit--explore is uniformly optimal; policy selection is frozen from validation utility--risk trade-offs.

\subsection{Source-level evidence diagnostics}
Let $\mathcal{B}_i$ be the one-class source annotations and $v(b,t)$ the intersection area between box $b$ and tile footprint $\Omega_t$, divided by the unclipped area of $b$. At visibility threshold $\eta$, source evidence recall is
\begin{equation}
R_{\mathrm{evid}}(i)=\frac{1}{|\mathcal{B}_i|}
\sum_{b\in\mathcal{B}_i}\mathbb{1}\!\left[
\max_{t\in\mathcal{S}_i}v(b,t)\geq\eta\right].
\label{eq:evidence-recall}
\end{equation}
Each source box receives credit at most once. Spatial coverage and overlap redundancy are
\begin{align}
C_{\mathrm{sp}}(i)&=\frac{|\bigcup_{t\in\mathcal{S}_i}\Omega_t|}{|\Omega_i|}, \\
D_{\mathrm{red}}(i)&=
\frac{\sum_{t\in\mathcal{S}_i}|\Omega_t|}
{|\bigcup_{t\in\mathcal{S}_i}\Omega_t|}-1.
\label{eq:spatial}
\end{align}
For empty sources, $R_{\mathrm{evid}}$ is undefined and excluded from the source-macro evidence-recall mean; negative-source behavior is reported separately. All footprints are clipped to $\Omega_i$, padding is excluded, and empty selected unions receive zero spatial coverage with redundancy marked undefined. Boundary truncation records covered instances with maximum visible fraction in $[\eta,.98)$. We report both negative retention (selected verified-negative candidates divided by all verified-negative candidates) and selected-negative fraction (selected verified-negative candidates divided by all selected candidates); neither is called recall without its denominator. These diagnostics measure available evidence, not task correctness.

\subsection{Structured multimodal inference}
Each selected tile and one frozen instruction are passed to a model-specific adapter. The current implementation uses prompt-only formatting followed by post-generation JSON extraction and validation; it is not grammar-constrained decoding. The frozen prompt, schema version, native image processor, generation parameters, raw output, and parser version are recorded. The parser records malformed output without silently repairing task fields. For generation $z$, constraint satisfaction and strict validity are distinct:
\begin{align}
s_{\mathrm{constr}}(z)&=\frac{1}{J}\sum_{j=1}^{J}\mathbb{1}[g_j(z)=1], \\
v_{\mathrm{schema}}(z)&=\prod_{j=1}^{J}\mathbb{1}[g_j(z)=1].
\label{eq:schema}
\end{align}
With $R>1$ generations, repeat consistency averages modal agreement over the four evaluated categorical fields $\mathcal{F}=\{\text{presence, count interval, density, annotated coverage}\}$:
\begin{equation}
s_{\mathrm{rep}}(t)=\frac{1}{|\mathcal{F}|}
\sum_{k\in\mathcal{F}}
\frac{\max_v\sum_{r=1}^{R}\mathbb{1}[z_{t,r,k}=v]}{R}.
\label{eq:consistency}
\end{equation}
Malformed or missing fields are retained as explicit disagreement. When $R=1$, repeat consistency is unavailable rather than set to one. The delivered tile prediction rule (first generation or frozen consensus) is declared before evaluation. Agreement is a stability diagnostic, not evidence of correctness. Detector agreement $s_{\mathrm{det}}$ compares VLM presence and count interval with fused detector evidence and is not treated as independent verification.

\subsection{Overlap-aware source reduction}
Detector boxes from all candidate tiles are projected to source coordinates and fused once using a validation-selected overlap threshold. Selected-tile VLM records annotate only their inspected footprints. Presence follows a frozen false-positive rule; count uses coordinate-aware fused instances; density uses count per valid source area; and annotated coverage uses the union of clipped fused boxes divided by valid source area. Categorical VLM fields use frozen reliability-weighted voting only where they are not deterministic functions of the fused boxes. Field conflicts and the uninspected-area fraction are retained as reducer features. The reducer returns a source record $\tilde y_i$ and reliability vector
\begin{equation}
u_i=[\bar v_{\mathrm{schema}},\bar s_{\mathrm{constr}},
\bar s_{\mathrm{rep}},\bar s_{\mathrm{det}},s_{\mathrm{conflict}},
s_{\mathrm{support}}].
\label{eq:sourcefeatures}
\end{equation}
If no valid VLM record remains, or if required evidence is unobserved under the frozen rule, the source is referred for review rather than relabeled empty. Reducer rules, overlap threshold, and field weights are frozen on validation sources.

\subsection{Validation-only selective decision}
A transparent fixed baseline averages the available reliability signals after orienting each so that larger is safer. The learned gate estimates source-record error probability $\hat r_i=h(u_i)$. To avoid fitting and thresholding on the same outcomes, source-grouped cross-fitting produces out-of-fold validation risks. For predeclared target risk $\epsilon$, the threshold is
\begin{equation}
\tau^*=\max_{\tau\in\mathcal{Q}_{\mathrm{val}}}\left\{\tau:
|A_\tau|\geq n_{\min},\quad
\frac{\sum_{i\in A_\tau}e_i}{|A_\tau|}\leq\epsilon\right\},
\label{eq:threshold}
\end{equation}
where $A_\tau=\{i:\hat r_i\leq\tau\}$, $e_i=1$ when any required source-level task field is incorrect, $\mathcal{Q}_{\mathrm{val}}$ is the frozen set of out-of-fold validation scores, and $n_{\min}$ is a predeclared minimum accepted count. If no threshold satisfies both constraints, the target is declared unattained; an empty accepted set is never reported as zero risk. Threshold ties use the highest-coverage rule followed by the lower-risk rule. After features and $\tau^*$ are frozen, the gate is refit on all validation sources and evaluated once on test sources. Records with $\hat r_i\leq\tau^*$ are accepted; all others are referred for review. This is empirical validation thresholding, not probability calibration or a distribution-free guarantee after crop shift.

\begin{algorithm}[t]
\caption{Fixed-budget source inference with selective review}
\label{alg:ornascout}
\begin{algorithmic}[1]
\Require $I_i$, tiler $(s,\alpha)$, detector/post-processing $f_D$, VLM $f_V$, prompt/schema $p$, repeats $R$, budget $K$, seed, frozen policy $\pi$, reducer $A$, gate $h_{\tau^*}$
\Ensure Source record $\tilde y_i$, decision $a_i$, tile audit trail $L_i$, and stage-cost ledger $G_i$
\State Initialize $L_i,G_i$; $\mathcal{T}_i\gets\mathcal{T}_{s,\alpha}(I_i)$ with coordinates and valid footprints
\ForAll{$t\in\mathcal{T}_i$}
  \State $\mathcal{D}_t\gets f_D(t)$; filter, clip, cache evidence, and update $G_i$
\EndFor
\State $\mathcal{S}_i\gets\pi(\mathcal{T}_i,\{\mathcal{D}_t\},\min(K,|\mathcal{T}_i|),\mathrm{seed}_i)$
\State $\mathcal{Z}_i\gets\emptyset$
\ForAll{$t\in\mathcal{S}_i$}
  \For{$r=1$ to $R$}
    \State Store raw $o_{t,r}\gets f_V(t,p)$; parse or mark malformed; update $L_i,G_i$
  \EndFor
  \State $\mathcal{Z}_i\gets\mathcal{Z}_i\cup\{(t,\Omega_t,o_{t,1:R})\}$
\EndFor
\State $(\tilde y_i,u_i)\gets A(\mathcal{Z}_i,\{(t,\Omega_t,\mathcal{D}_t):t\in\mathcal{T}_i\})$
\State $a_i\gets\mathrm{accept}$ if $h(u_i)\leq\tau^*$, else $\mathrm{review}$
\State \Return $\tilde y_i,a_i,L_i,G_i$
\end{algorithmic}
\end{algorithm}
